%=============================================================================
%  08_threats.tex -- Threats to validity
%
%  CONVENTION: as in 06_applications.tex.
%=============================================================================
\section{Threats to Validity}
\label{sec:threats}

We list every known threat, including the ones that weaken results we would
prefer to keep. They are ordered by how much damage they do.

\paragraph{1. The comparison is not size-matched, and this is the most serious
threat.} Every diffusion measurement uses \modelLLaDA{} or \modelDG{} against
\modelARthree. The diffusion models are larger --- substantially so by total
parameter count, and \modelDG{} is a mixture-of-experts model whose active
parameter count is closer to the AR baseline than its total suggests. This does
\emph{not} affect \cref{sec:cost-laws,sec:two-walls,sec:mechanism}: a cost law is a
per-model measurement, and the batching multiplier is a property of each model's
own scaling, so those results stand as measured regardless of parameter parity.
It \emph{does} affect \cref{sec:applications} in full. Every application-level
quality comparison, and every latency ratio interpreted as an architectural
difference rather than a size difference, is confounded. The size-matched arm
(\modelDG{} against \modelARseven) is filed as issue~\#34 and is not yet run.
\open{\#34}

\paragraph{2. The AR cost-law fit is not reconciled across artifacts.} The
fitted intercept, slope and residual $\sigma$ we report come from the analysis
JSON; an earlier report generated from the same campaign records a different
intercept and a residual $\sigma$ nearly an order of magnitude smaller, and an
internal digest records a third variant with no intercept at all. We use the
JSON, but the discrepancy is unexplained. It must be resolved by re-fitting from
the raw per-cell data before submission; the slope --- which carries the
marginal-cost claim --- agrees across all three, but the intercept and the
uncertainty estimate do not. \open{fit reconciliation}

\paragraph{3. The two models are served asymmetrically, and we did not choose
this.} \modelDG{} runs at \texttt{--max-num-seqs 1} while the AR baseline runs
the engine's default 256-way scheduler. The asymmetry is forced by an engine
crash that fires as soon as two requests in different denoising phases are
scheduled together (\cref{sec:two-walls}); serialised serving is therefore
diffusion's production-relevant mode on this engine today, but it is not a
choice we made in the experiment's favour or against it, and a fixed engine
would change the diffusion numbers and not the AR ones.

\paragraph{4. All dInfer latencies are cache-less.} dInfer's prefix and vicinity
KV-cache paths error on prompts at or above \dinferCachePromptLimit{} tokens,
which our prompts exceed, so every dInfer latency we report was measured with
caching disabled. dInfer's optimised speed on this task is therefore unmeasured.
This is a limitation of that runtime at our prompt length, not a property of
diffusion, and the \cref{sec:neg-adaptive} latency figures in particular should
be read as an upper bound rather than a measurement of adaptive decoding.

\paragraph{5. AR's record completeness may be a harness artifact.} As set out in
\cref{sec:app-deadline}, AR reports \arRecordComplete\% record completeness at
every deadline including ones it finishes well inside, which is more consistent
with a parsing or prompt-preamble defect than with truncation. Issue~\#38 is
open. That result should not be cited as an architectural finding until it
closes. \open{\#38}

\paragraph{6. A single corpus.} Every number here is BGL. The paper whose
claims motivated this work averages over four LOGHUB corpora; we do not. Log
corpora differ substantially in template diversity and line structure, and we
cannot rule out that some of what we attribute to architecture is a property of
BGL's format.

\paragraph{7. Poisson arrivals understate log burstiness.} The open-loop driver
in \cref{sec:two-walls} samples inter-arrival times from an exponential
distribution. Real log streams are burstier, and a bursty rerun would move both
saturation knees to lower mean rates. Latency past capacity is measured in
transient windows with a bounded drain, so the location of the knee is the
result we defend; the exact tail milliseconds beyond it are not.

\paragraph{8. Semantic gold labels are self-generated.} No teacher model was
released with the work we are validating, so the gold labels for the three
semantic fields are placeholders we generated. This caps achievable semantic F1
for every model equally and does not bias the comparison, but it does mean the
absolute Field-F1 numbers in this paper are not comparable to numbers produced
against a different label source.

\paragraph{9. We never ran the model we are comparing against.} FlowFill's
weights are unreleased. Every comparison to it in this paper is a comparison to
its published claims, not to a run of it. We have not reproduced its numbers,
have not failed to reproduce them, and make no statement about whether they are
attainable --- only about what our own substitutes do under our own harness.

\paragraph{10. Small $n$ in places.} The semi-autoregressive block ablation
(\cref{sec:neg-semiar}) is $n=30$ and is labelled preliminary throughout. The
adaptive-effort experiment (\cref{sec:neg-adaptive}) uses 100 events per
difficulty bin for diffusion against 200 per bin for AR. Neither supports a
conclusion at the precision at which the numbers are printed.
